%-------------------------
% Resume in Latex
% Author : Vidushi Wahal
%------------------------

\documentclass[letterpaper,11pt]{article}

    \newcommand{\updateinfo}[1][\today]{\par\vfill\hfill{\scriptsize\color{black}Last updated on #1}}


\usepackage{changepage}
\usepackage{latexsym}
\usepackage[empty]{fullpage}
\usepackage{titlesec}
\usepackage{marvosym}
\usepackage[usenames,dvipsnames]{color}
\usepackage{verbatim}
\usepackage{enumitem}
\usepackage[pdftex]{hyperref}
\usepackage{fancyhdr}
\usepackage{ragged2e}
\usepackage{tabularx}


\pagestyle{fancy}
\fancyhf{} % clear all header and footer fields
\fancyfoot{}
\renewcommand{\headrulewidth}{0pt}
\renewcommand{\footrulewidth}{0pt}

% Adjust margins
\addtolength{\oddsidemargin}{-0.375in}
\addtolength{\evensidemargin}{-0.375in}
\addtolength{\textwidth}{1in}
\addtolength{\topmargin}{-.5in}
\addtolength{\textheight}{1.0in}

\urlstyle{same}

\raggedbottom
\raggedright
\setlength{\tabcolsep}{0in}

% Sections formatting
%\titleformat{\section}{
%  \vspace{-4pt}\scshape\raggedright\large
%}{}{0em}{}[\color{black}\titlerule \vspace{-5pt}]

%-------------------------
% Custom commands
%\newcommand{\resumeItem}[2]{
%  \item\small{
%    \textbf{#1}{: #2 \vspace{-2pt}}
%  }
%}
%
%\newcommand{\resumeSubheading}[4]{
%  \vspace{-1pt}\item
%    \begin{tabular*}{0.97\textwidth}{l@{\extracolsep{\fill}}r}
%      \textbf{#1} & #2 \\
%      \textit{\small#3} & \textit{\small #4} \\
%    \end{tabular*}\vspace{-5pt}
%}
%
%\newcommand{\resumeSubItem}[2]{\resumeItem{#1}{#2}\vspace{-4pt}}
%
%\renewcommand{\labelitemii}{$\circ$}
%
%\newcommand{\resumeSubHeadingListStart}{\begin{itemize}[leftmargin=*]}
%\newcommand{\resumeSubHeadingListEnd}{\end{itemize}}
%\newcommand{\resumeItemListStart}{\begin{itemize}}
%\newcommand{\resumeItemListEnd}{\end{itemize}\vspace{-5pt}}

%-------------------------------------------

%%%%%% REFRENCES SECTION

%%%%%%  CV STARTS HERE  %%%%%%%%%%%%%%%%%%%%%%%%%%%%


\begin{document}

%----------HEADING-----------------
\begin{tabular*}{\textwidth}{l@{\extracolsep{\fill}}r}
  \textbf{{\Large Michael Fabian Barczay}} & \\
   & \\
\end{tabular*}
\begin{tabular*}{\textwidth}{l@{\extracolsep{\fill}}r}
  European University Institute & \textbf{Email}: \href{mailto:michael.barczay@eui.eu}{michael.barczay@eui.eu}\\
  Department of Economics &\textbf{Website}: \href{https://michaelbarczay.com/}{michaelbarczay.com}  \\
  Via delle Fontanelle 18 & \textbf{Phone}: [+41] 079 862 18 66 \\
  50014 Fiesole, Italy &  
\end{tabular*}

%-----------Research Interests-----------------
\vspace{1cm}
\textbf{RESEARCH INTERESTS}\\
\vspace{1mm} \hrule
\vspace{2mm}
\hspace{1.2cm} \textbf{Primary Fields:} (Quantitative) Macroeconomics, Public Finance\\
\hspace{1.2cm} \textbf{Interests:} Heterogeneous Agents, Optimal Taxation, Consumption, Green Transition\\

%-----------Research Interests-----------------
\vspace{0.5cm}
\textbf{CURRENT POSITIONS}\\
\vspace{1mm} \hrule
\vspace{2mm}
\hspace{1.2cm} \textbf{PhD Candidate,} European University Institute (Department of Economics), Italy\\
\hspace{2.1cm} Supervisors: Alexander Monge-Naranjo, Russell Cooper\\
\hspace{1.2cm} \textbf{Research Associate,} Study Center Gerzensee, Switzerland

%-----------Education-----------------
\vspace{0.5cm}
\textbf{EDUCATION}\\
\vspace{1mm} \hrule
\vspace{2mm}
\hspace{0.9cm} \textbf{PhD in Economics}, European University Institute \hfill 2020 - 2025\\ \vspace{1mm}
\hspace{0.9cm} \textbf{MRes in Economics}, European University Institute \hfill 2020 - 2021\\ \vspace{1mm}
%\hspace{1.3cm} PhD courses in Microeconomics, Econometrics and Macroeconomics \hfill 2020-2021\\ \vspace{3mm}
\hspace{0.9cm} \textbf{MA Specialized Economic Analysis}, Barcelona School of Economics \hfill{2018 - 2019}\\  \vspace{1mm}
%\hspace{1.3cm} Program in International Trade, Finance and Development \hfill 2018 - 2019\\ \vspace{0mm}
%\hspace{1.3cm} Thesis Supervisors: Fernando Broner, Antonio Ciccone\\ \vspace{3mm}
\hspace{0.9cm} \textbf{MA International Affairs}, Geneva Graduate Institute \hfill{2016 - 2018}\\ \vspace{1mm}
%\hspace{1.3cm}  \hfill 2016 - 2018\\ \vspace{3mm}
%\hspace{0.9cm} Exchange Semester, American University in Cairo\hfill{\textbf{Cairo, Egypt}}\\ \vspace{1mm}
%\hspace{1.3cm} Courses in Economics and Political Science  \hfill Autumn 2017\\ \vspace{3mm}
\hspace{0.9cm} \textbf{BA Business and Economics} (Major in Economics), University of Basel \hfill{2013 - 2016}\\ 
%\hspace{1.3cm} Major in Economics\hfill 2013 - 2016\\ \vspace{0mm}
%\hspace{1.3cm} Thesis Supervisor: Sarah M. Lein
%-----------PHD MACRO COURSES-----------------

%\vspace{0.5cm}
%\textbf{MACROECONOMICS PHD COURSES}\\
%\vspace{2mm}
%\hspace{0.9cm} \textbf{Core Courses Year 1}  \\ \vspace{0mm}
%\begin{adjustwidth}{1.3cm}{0cm} Macroeconomics I (R. Cooper; OLG with Fiscal and Monetary Policy), Macroeconomics II (J. Bueren and E. Challe; RBC, Dynamic Programming, New Keynesian Model), Macroeconomics III (A. Monge-Naranjo and E. Challe; Incomplete Markets, Search \& Matching)\\
%\end{adjustwidth}	
%
%
%\vspace{2mm}
%\hspace{0.9cm} \textbf{Advanced Courses Year 1/2}\\ \vspace{0mm}
%\begin{adjustwidth}{1.3cm}{0cm}
%Endogenous and Exogenous Incomplete Markets (\'A. \'Abrah\'am), Fiscal and Monetary Policy (R. Marimon), Life-Cycle Heterogenous Agent Models (J. Bueren), Computations and Quantitative Models in Macro (A. Monge-Naranjo), Optimal Fiscal Policy in Quantitative Macro Models (A. Ferriere), Inequality and Education (A. Monge-Naranjo), Firm Heterogeneity (B. Grassi), Financial Frictions (S. Tsiaras)
%\end{adjustwidth}


%-----------Experience-----------------
\vspace{0.5cm}
\textbf{EXPERIENCE}\\
\vspace{1mm} \hrule
\vspace{2mm}
\hspace{0.9cm} \textbf{Fund Intern}, International Monetary Fund (Fiscal Affairs Department) \hfill 06/2024 - 08/2024\\ \vspace{1mm}
%\hspace{0.9cm} Data Replicator, Journal of the European Economic Association \hfill 02/2024 - today\\ \vspace{0mm}
%\vspace{2mm}
\hspace{0.9cm} \textbf{Visiting PhD Student}, Banco Central de Chile (Research Department) \hfill 08/2023 - 10/2023\\ \vspace{1mm}
\hspace{0.9cm} \textbf{Statistical Assistant}, European University Institute \hfill 04/2023 - 09/2024\\ \vspace{1mm}
\hspace{0.9cm} \textbf{Research Assistant} to Russell Cooper, European University Institute \hfill 03/2023 - 05/2023\\ \vspace{1mm}
\hspace{0.9cm} \textbf{Teaching Assistant \& Lecturer}, European University Institute \hfill 09/2021 - 03/2024\\ \vspace{1mm}
%\hspace{0.9cm} \textbf{Research Assistant} to Alexander Monge-Naranjo, European University Institute \hfill 12/2022\\ \vspace{1mm}
\hspace{0.9cm} \textbf{Intern}, Swiss National Bank (Inflation Forecasting Unit) \hfill 08/2019 - 07/2020 \\ \vspace{1mm}
\hspace{0.9cm} \textbf{Research Assistant} to David Sylvan, Geneva Graduate Institute \hfill 01/2018-07/2018 \\ 

%-----------Teaching-----------------
\vspace{0.5cm}
\textbf{TEACHING}\\
\vspace{1mm} \hrule
\vspace{2mm}
\hspace{0.9cm} Tax Systems\hfill Fall 2024\\ 
\hspace{1.3cm} SC Gerzensee, PhD Level -- Teaching Assistant to Joel Slemrod (U. Michigan)\\ 
\vspace{1mm}
\hspace{0.9cm} Life-Cycle Heterogeneous Agents Models: Solution and Estimation\hfill Fall 2022 \& Spring 2024\\ 
\hspace{1.3cm} EUI, PhD Level -- Teaching Assistant to Jes\'us Bueren\\ 
\vspace{1mm}
\hspace{0.9cm} Matlab Programming Bootcamp\hfill Fall 2021, 2022 \& 2023\\ \vspace{0mm}
\hspace{1.3cm} EUI, PhD Level -- Lecturer\\ 
\vspace{1mm}
\hspace{0.9cm} Macroeconomics I (Dynamic Fiscal and Monetary Policy) \hfill Fall 2021\\ \vspace{0mm}
\hspace{1.3cm} EUI, PhD Level -- Teaching Assistant to Russell Cooper\\



\pagebreak
%-----------WORK-----------------
\vspace{0.5cm}
\textbf{WORK IN PROGRESS}\\
\vspace{1mm} \hrule
\vspace{2mm}
\begin{adjustwidth}{0.9cm}{0cm} ``On the Optimal Design of Consumption Taxes'' \textit{(EUI Best Second Year Paper Award)}\end{adjustwidth} 
\begin{adjustwidth}{1.1cm}{1.4cm}\footnotesize \justify{\textit{Abstract:} How should differentiated consumption taxes be designed in the presence of capital income taxes and progressive labor income taxes?
I study this question using a quantitative model featuring heterogeneous households with non-homothetic preferences, uninsurable idiosyncratic risk, and a government that uses various tax instruments to raise revenue. I estimate the parameters governing households' demand using data from the US Consumption Expenditure Survey, and show that my model matches the heterogeneous consumption behavior across the income distribution. Allowing the benevolent government to jointly optimize consumption taxes on 11 different consumption categories and labor income taxes, I find that necessities should be heavily subsidized (-52\%) and that luxuries are optimally taxed at a positive rate of 7\%. This increase in the rate differential is optimally compensated by a significant decrease in the progressivity of the labor income tax. Three main mechanisms explain why such differentiated tax rates are welfare maximizing: they provide consumption insurance by subsidizing essential goods of low-income households, imply a targeted taxation of the initial wealth of high-wealth households, and induce highly productive households to increase their labor supply.}\end{adjustwidth} 

 \vspace{0.1cm}
 \begin{adjustwidth}{0.9cm}{1.4cm} ``How to Finance the Green Transition: the Political Economy of Investment Tax Credits'' (with \textit{Russell Cooper})\end{adjustwidth} 
 
 \begin{adjustwidth}{1.1cm}{1.4cm}\footnotesize \justify{\textit{Abstract:} We study the aggregate and distributional consequences of green investment tax credits (ITCs) and ask under which financing structures such environmental policies would be adopted by a majority of voters and sustained in the long run.
              We develop an overlapping generations model with heterogeneous households, multiple sectors, and a government that wants to introduce a green ITC to reduce pollution. Our model highlights both an intratemporal (across the income distribution) and an intertemporal (across generations) disagreement about the desirability of green ITCs arising from the unequal distribution of the costs and benefits. Together, they can lead to voting outcomes in which the ITC would never be adopted, even though it would be welfare improving for a majority of the population in the long run. We show that allowing for some debt financing of the ITC can overcome this political gridlock. Moreover, this debt can be fully repaid in the long run while maintaining high approval rates for the ITC. Changes in asset market participation rates and factor prices induced by the ITC explain why fully tax-financed ITCs are approved only in the long run, but not at the time of introduction of the ITC.}\end{adjustwidth}
 
  \vspace{0.1cm}
\begin{adjustwidth}{0.9cm}{1.4cm} ``Taxing Mobile Money: Theory and Evidence'' (with \textit{Shafik Hebous, Fay\c{c}al Sawadogo, Jean-Francois Wen)}\end{adjustwidth} 

 \begin{adjustwidth}{1.1cm}{1.4cm}\footnotesize \justify{\textit{Abstract:} Mobile money has emerged as a widespread digital alternative to traditional banking, transforming financial services and inclusion in many developing countries. We develop a stylized model of cash management to analyze the effects of mobile money taxation. The model predicts that such a tax incentivizes users to switch to alternative payment systems, disproportionately affecting individuals with limited banking access. We empirically test these predictions using two different types of data. First, we use cross-country financial access surveys, leveraging the staggered introduction of mobile money taxes in Sub-Saharan Africa. We find that the tax reduces the number of mobile money users and the frequency and value of transactions. Second, we use mobile money transaction-level data from Cameroon and Mali. We find that the mobile money tax reduces the probability of using mobile money by about 10 percentage points. The total transaction volume is reduced by about 40\%, suggesting a very high tax elasticity. Unbanked individuals and rural areas respond less strongly to the introduction of the tax, suggesting that these users pay a disproportionate share of the tax.}\end{adjustwidth}

 
   \vspace{0.1cm}
\begin{adjustwidth}{0.9cm}{1.4cm} ``The Regional Incidence of VAT Tax Cuts: Evidence from Germany'' (with \textit{Shafik Hebous, Tom Zimmermann)}\end{adjustwidth} 
 
 \vspace{0.1cm}
\begin{adjustwidth}{0.9cm}{1.4cm} ``Heterogeneity in the Bank Lending Channel: Evidence from the Euro Area'' (with \textit{Sylvia Kaufmann})\end{adjustwidth} 

 \vspace{0.1cm}
\begin{adjustwidth}{0.9cm}{1.4cm} ``Participation and Capital Dynamics: A Framework for Analysis'' (with \textit{Russell Cooper})\end{adjustwidth} 

%\begin{adjustwidth}{1.1cm}{1.4cm}\footnotesize \justify{\textit{Abstract:} In this paper we develop a novel Bayesian framework to estimate heterogeneous and time-varying responses of agents to monetary policy shocks. Monte Carlo simulations illustrate the performance of the method. We then apply the framework to euro area data to study the heterogeneous transmission of the ECB’s monetary policy across euro area countries. Our results suggest that the bank lending has been operating heterogeneously across euro area countries and with substantial heterogeneity across the last 25 years.}\end{adjustwidth}

% \vspace{0.1cm}
%\begin{adjustwidth}{0.9cm}{1.4cm} ``The Anatomy of the Marginal Propensity to Consume'' (with \textit{Lucciano Villacorta, waiting for data})\end{adjustwidth} 

%\begin{adjustwidth}{1.1cm}{1.4cm}\footnotesize \justify{\textit{Abstract:} In this project, we aim to make use of administrative data from Chile to better understand how income shocks translate into consumption responses. Specifically, we ask whether households respond to unexpected income changes by adjusting the level of consumption, the composition of the consumption basket in terms of quality composition and expenditure categories, or by searching for lower prices. For this project, we plan to extract income shocks from income tax data and combine it with detailed transaction level data.}\end{adjustwidth}

%-----------CONFERENCES-----------------
\vspace{0.5cm}
\textbf{CONFERENCES \& SEMINARS} \\
\vspace{1mm} \hrule
\vspace{1mm}
\begin{adjustwidth}{0.9cm}{0cm} Workshop on ``Digital Money and Taxation'', \textit{Goethe University, Frankfurt}, 7.3.2025 \tiny \normalsize \end{adjustwidth}
\begin{adjustwidth}{0.9cm}{0cm} Econometric Society Winter Meeting, \textit{The University of the Balearic Islands}, 16. - 18.12.2024 \tiny \normalsize \end{adjustwidth}
\begin{adjustwidth}{0.9cm}{0cm} SED Winter Meeting, \textit{Universidad Torcuato di Tella}, 12. - 14.12.2024 \tiny \normalsize\end{adjustwidth}
\begin{adjustwidth}{0.9cm}{0cm} Doctoral Workshop on Quantitative Dynamic Economics, \textit{University of Aix-Marseille}, 14.10.2024\end{adjustwidth}
\begin{adjustwidth}{0.9cm}{0cm} MMF Annual Conference '24, \textit{University of Manchester}, 06.09.2024\end{adjustwidth}
\begin{adjustwidth}{0.9cm}{0cm} Tax Policy Brown Bag , \textit{IMF}, 28.08.2024\end{adjustwidth}
\begin{adjustwidth}{0.9cm}{0cm} \textit{Paris School of Economics}, 25.04.2024\end{adjustwidth}
\begin{adjustwidth}{0.9cm}{0cm} \textit{Study Center Gerzensee}, 25.04.2024\end{adjustwidth}
\begin{adjustwidth}{0.9cm}{0cm} Simposio of the Spanish Economic Association (SAEe) \textit{University of Salamanca}, 16.12.23\end{adjustwidth}
%\vspace{1mm}
\begin{adjustwidth}{0.9cm}{0cm} \textit{Banco Central de Chile}, 17.10.23\end{adjustwidth}
%\vspace{1mm}
\begin{adjustwidth}{0.9cm}{0cm} 17th End-of-Year Conference of Swiss Economists Abroad, \textit{University of Basel}, 22.12.22\end{adjustwidth}

\vspace{0.5cm}
\textbf{SHORT COURSES} \\
\vspace{1mm} \hrule
\vspace{1mm}
\begin{adjustwidth}{0.9cm}{0cm} Bayesian Methods for Empirical Macroeconomics, Gary Koop, \textit{Studycenter Gerzensee}, July 2022\end{adjustwidth}
\begin{adjustwidth}{0.9cm}{0cm} Tools for Macroeconomists (Advanced), Wouter den Haan \& Pontus Rendahl, \textit{Oxford}, August 2021\end{adjustwidth}
\begin{adjustwidth}{0.9cm}{0cm} Inference in Structural Macro Models, Domenico Giannone \& Girogio Primiceri, \textit{Swiss National Bank}, January 2020\end{adjustwidth}



%-----------SKILLS-----------------
\vspace{0.5cm}
\textbf{SCHOLARSHIPS \& AWARDS}\\
\vspace{1mm} \hrule
\vspace{2mm}
\hspace{0.9cm} SAEe PhD Student Grant, Fundacion Ramon Areces\hfill 2023\\
\vspace{1mm}
\hspace{0.9cm} EUI Best Second Year Paper Award\hfill 2022\\
\vspace{1mm}
\hspace{0.9cm} EUI Scholarship Swiss State Secretariat for Education, Research and Innovation\hfill since 2020\\
\vspace{1mm}
\hspace{0.9cm} Dr. Max Huusmann Foundation, Zurich, Scholarship for the Graduate Institute \hfill 2016-2018\\

%-----------SKILLS-----------------
\vspace{0.5cm}
\textbf{SKILLS}\\
\vspace{1mm} \hrule
\vspace{1mm}
\begin{adjustwidth}{0.9cm}{0cm}\textbf{Computational}\end{adjustwidth}
\begin{adjustwidth}{1.3cm}{0cm} Julia, Matlab, Stata, R \end{adjustwidth}
\vspace{0.2cm}
\begin{adjustwidth}{0.9cm}{0cm}\textbf{Languages}\end{adjustwidth} 
\begin{adjustwidth}{1.3cm}{0cm} German (native), English (fluent), French (proficient), Italian (proficient), Spanish (basic), Arabic (beginner)\end{adjustwidth}


%-----------SKILLS-----------------
\vspace{0.5cm}
\textbf{PROFESSIONAL ACTIVITIES, REFEREEING \& VOLUNTEERING}\\
\vspace{1mm} \hrule
\vspace{0.2cm}
\hspace{1.3cm} Referee for: International Tax and Public Finance\\
\vspace{0.1cm}
\hspace{1.3cm} Matlab Tutor, European University Institute \hfill 2022-2023\\
\vspace{0.1cm}
\hspace{1.3cm} Co-Organizer, Macro Working Group, European University Institute \hfill 2021-2022\\
\vspace{0.1cm}
\hspace{1.3cm} Researchers Representative Economics Department, European University Institute \hfill 2021\\ \vspace{0.1cm}
\hspace{1.3cm} Treasurer \& Board Member, Pro Centro Educativo Yampuc \hfill since 2016\\
\vspace{0.1cm}
%\hspace{1.3cm} Ambassador, European University Institute \hfill since 2021\\

%\pagebreak
%-----------REFERENCES-----------------
%\vspace{0.5cm}
%\textbf{REFERENCES}\\
%\vspace{1mm} \hrule
%\vspace{0.25cm}
%\begin{tabularx}{\textwidth}{XXX}
%\textbf{Alexander Monge-Naranjo}  &  \textbf{Russell Cooper}& \textbf{\'Arp\'ad \'Abrah\'am}  \\
%European University Institute & European University Institute & University of Bristol  \\
%Department of Economics & Department of Economics & Department of Economics  \\
%Via delle Fontanelle 18 & Via delle Fontanelle 18 & Priory Road Complex, Priory Road  \\
%50014 Fiesole, Italy & 50014 Fiesole, Italy & BS81TU, United Kingdom \\
%\href{mailto:alexander.monge-naranjo@eui.eu}{alexander.monge-naranjo@eui.eu}& \href{mailto:russell.cooper@eui.eu}{russell.cooper@eui.eu}  & \href{mailto:arpad.abraham@bristol.ac.uk}{arpad.abraham@bristol.ac.uk} \\\
%[+39] 055 4685 942 & [+39] 055 4685 908 & [+44] 117 4553 064 \\
%
%\end{tabularx}

%\begin{center}
%\hspace{-4.2cm}	
%\begin{tabular}{lll}
%\textbf{Alexander Monge-Naranjo} &\hspace{2.2cm} &\textbf{Russell Cooper} \hspace{2.6cm} \textbf{Arpad Abraham} \\
%European University Institute            &\hspace{1cm}& European University Institute               \\
%Department of Economics           &\hspace{2.2cm}& Department of Economics                \\
%Via delle Fontanelle 18
%  &\hspace{2.2cm}& Via delle Fontanelle 18 \\
%50014 Fiesole, Italy
% &\hspace{2.2cm}& 50014 Fiesole, Italy    \\
% \url{alexander.monge-naranjo@eui.eu}       & \hspace{2.2cm}& \url{russell.cooper@eui.eu}\\
%
%  [+39] 055 4685 942       & \hspace{2.2cm} & [+39] 055 4685 908                            
%\end{tabular}
%\end{center}


\updateinfo

\end{document}